\documentclass{beamer}
\usetheme{metropolis} % Use the metropolis theme




% Add tikz and pgfplots packages
\usepackage{tikz, pgfplots}
\usetikzlibrary{positioning}

% For clicking references
\usepackage{hyperref}

% For better referencing
\usepackage{cleveref}

\usepackage{graphicx}

% Define custom pastel colors
\definecolor{pastelRed}{RGB}{255, 105, 97}   % A soft pastel red
\definecolor{pastelBlue}{RGB}{119, 158, 203} % A muted pastel blue
\definecolor{pastelYellow}{RGB}{255, 223, 0} % A gentle pastel yellow
\definecolor{lightGray}{RGB}{211, 211, 211}  % A light gray for subtitles and less emphasized text

% Apply the custom colors
\setbeamercolor{palette primary}{bg=black, fg=white}
\setbeamercolor{palette secondary}{bg=lightGray, fg=black}
\setbeamercolor{palette tertiary}{bg=black, fg=white}
\setbeamercolor{titlelike}{parent=palette primary, fg=black}
\setbeamercolor{subtitle}{fg=lightGray}
\setbeamercolor{structure}{fg=black} % For itemize, enumerate, etc

% Change color of normal text
\setbeamercolor{normal text}{fg=black, bg=white}

% Set the color of the table of contents
\setbeamercolor{section in toc}{fg=black} % Section titles in TOC
\setbeamercolor{subsection in toc}{fg=black} % Subsection titles in TOC

% Set block colors
\setbeamercolor{block title}{use=structure,fg=white,bg=pastelRed}
\setbeamercolor{block body}{fg=black,bg=white}



% Title Page Info
\title{Discrete Probability, Random Variables, and Bounds}
\subtitle{Spørgsmål 3 fra Exam Questions}
\author{Kevin Vinther}
\date{\today}

\begin{document}

% Title Page
\begin{frame}
  \titlepage
\end{frame}

% Table of Contents
\begin{frame}[allowframebreaks]
  \frametitle{Table of Contents}
  \tableofcontents
\end{frame}

\section{An Introduction to Discrete Probability}

\begin{frame}[allowframebreaks]
  \frametitle{Finite Probability}
  \begin{itemize}
  \item \textbf{Experiment}: En procedure som giver et resultat ud af et givet sæt a mulige resultater.
  \item \textbf{Sample Space}: Sættet af alle mulige resultater.
  \item \textbf{Event}: Subset af sample space.
  \end{itemize} 
\end{frame}

\begin{frame}
  \frametitle{Definition 1}
  \begin{definition}
    If $S$ is a finite nonempty sample space of equally likely outcomes, and $E$ is an event, that is, a subset of $S$, then the \textit{probability} of $E$ is $p(E) = \frac{|E|}{|S|}$.
  \end{definition}
  \begin{itemize}
  \item \textit{Implikationer:}
  \item Givet, at $E$ er en event, og $S$ sample space:
  \item $0 \leq |E| \leq |S|$, fordi $E \subseteq S$. Dermed $0 \leq p(E) = |E|/|S| \leq 1$.
  \end{itemize}
\end{frame}

\begin{frame}
  \frametitle{Example 1}
  \begin{itemize}
  \item<1-> En boks indeholder fire blå bolde og fem røde. Hvad er sandsynligheden at en bold der er valgt tilfældigt er blå?
  \item<2-> $\frac{4}{9}$
  \end{itemize}
\end{frame}

\begin{frame}
  \frametitle{Example 2}
  \begin{itemize}
  \item<1-> Hvad er sandsynligheden for at, når 2 terninger bliver rullet, at summen af de to tal på terningen er 7?
  \item<1-> $6^2 = 36$ mulige udfald. (Af product rule).
  \item<2-> Der er 6 udfald der ender i 7: (1,6), (2,5), (3,4), (4,3), (5,2), og (6,1)
  \item<3-> Dermed: $\frac{6}{36} = \frac{1}{6}$
  \end{itemize} 
\end{frame}

\begin{frame}
  \frametitle{Eksempel 5}
  \begin{itemize}
  \item<1-> Hvad er sandsynligheden for at en hånd i poker har 4 kort af én slags (samme tal)? (Hint: Product rule)
  \item<2-> Antallet af hænder med 5 kort med fire af én slags er produktet af antallet måder man kan vælge én slags på, antallet af måder du kan vælge disse 4 af slagsen ud af de fire i dækket, og antallet af måder du kan vælge det femte kort. Det er:
  \item<3-> $\binom{13}{1}\binom{4}{4}\binom{48}{1}$. Næste: Hvor mange forskellige hænder af 5 kort er der? 
  \item<4-> $\binom{52}{5}$. Næste: Find sandsynligheden
  \item<5-> $\frac{\binom{13}{1}\binom{4}{4}\binom{48}{1}}{\binom{52}{5}}= \frac{13 \cdot 1 \cdot 48}{2598960} \approx 0.00024$
  \end{itemize}  
\end{frame}

\begin{frame}
  \frametitle{Eksempel 7 i}
  \begin{itemize}
  \item<1-> Hvad er sandsynligheden at tallene 11, 4, 17, 39 og 23 bliver taget i den orden fra en spand med 50 bolde med tallene 1..50, hvis (a) bolden der er valgt ikke bliver lagt tilbage i spanden før den næste bold blive valgt og (b) bolden bliver lagt tilbage før den næste bold bliver valgt?
  \item<1-> Lad os starte med spørgsmål (a) (hint: Hvor mange måder kan man vælge boldene på?):
  \item<2-> Man kan vælge 5 bolde på $50 \cdot 49 \cdot 48 \cdot 47 \cdot 46 = 254251200$ måder. Hvad er sandsynligheden så for at vælge præcis den orden der er beskrevet tidligere?
  \item<3-> $\frac{1}{254251200}$.
  \item<3-> Bogen kalder dette \textbf{sampling without replacement}.
  \end{itemize} 
\end{frame}

\begin{frame}
  \frametitle{Eksempel 7 ii}
  \begin{itemize}
  \item<1-> Hvad med (b)? Hvor mange måder kan man vælge 5 bolde på så?
  \item<2-> $50^5$. Husk at antallet af permutationer med gentagelser er $n^r$, og her er $n = 50, r = 5$.
  \item<3-> Måden man finder sandsynligheden på har ikke ændret sig her.
  \item<3-> $\frac{1}{312500000}$
  \item<3-> Bogen kalder dette \textbf{sampling with replacement}.
  \item<3-> Hvad forskellen er mellem det og permutations with repetition er et godt spørgsmål.
  \end{itemize} 
\end{frame}

\subsection{Probabilities of Complements and Unions of Events}
\label{subsec:probcompunion}

\begin{frame}
  \frametitle{Theorem 1}
  Tid til første theorem!
  \begin{theorem}[Theorem 1]
    Let $E$ be an event in a sample space $S$. The probability of the event $\overline{E} = S - E$, the complementary event of $E$ is given by $$P(\overline{E}) = 1 - p(E)$$
  \end{theorem}
\end{frame}

\begin{frame}
  \frametitle{Theorem 1 Bevis}
  \begin{proof}[Theorem 1 Bevis]
    For at finde sandsynligheden af begivenheden $\overline{E} = S - E$, notér at $|\overline{E}| = |S| - |E|$. Dermed: $$p(\overline{E}) = \frac{|S|-|E|}{|S|}=1 - \frac{|E|}{|S|}=1-p(E)$$
  \end{proof}
\end{frame}

\begin{frame}
  \frametitle{Eksempel 8}
  \begin{itemize}
  \item<1-> En sekvens af 10 bits er tilfældigt genereret. Hvad er sandsynligheden for at mindst en af disse bits er 0?
  \item<1-> Hint: Brug complement.
  \item<2-> Lad $E$ være tilfældet at mindst en af de 10 bits er 0. Så er $\overline{E}$ sandsynligheden for at alle bitsne er 1. 
  \item<3-> $P(E) = 1-p(\overline{E})=1- \frac{|\overline{E}|}{|S|}=1- \frac{1}{2^{10}} = 1 - \frac{1}{1024} = \frac{1023}{1024}$
  \end{itemize}
\end{frame}

\begin{frame}
  \frametitle{Theorem 2}
  \begin{theorem}
    Let $E_{1}$ and $E_{2}$ be events in the sample space $S$. Then $$p(E_{1} \cup E_{2}) = p(E_{1})+P(E_{2})-p(E_{1} \cap E_{2}).$$
  \end{theorem} 
\end{frame}

\begin{frame}
  \frametitle{Theorem 2 Bevis}
  \begin{proof}[Theorem 2 Bevis]
    Vi bruger subtraction rule for sæt (inclusion-exclusion for 2 sæt):
    $$|E_{1} \cup E_{2}| = |E_{1}| + |E_{2}| - |E_{1} \cap E_{2}|$$
    Dermed:
    \begin{equation*}
      \begin{split}
        p(E_{1} \cup E_{2}) &= \frac{|E_{1} \cup E_{2}|}{|S|}\\
                         &= \frac{|E_{1}|+|E_{2}-|E_{1}\cap E_{2}|}{|S|}\\
                         &= \frac{|E_{1}|}{|S|} + \frac{|E_{2}|}{|S|}- \frac{|E_{1} \cap E_{2}|}{|S|}\\
                         &= p(E_{1})+p(E_{2})-p(E_{1} \cap E_{2})
      \end{split}
    \end{equation*}
  \end{proof}
\end{frame}

\begin{frame}
  \frametitle{Eksempel 9}
  \begin{itemize}
  \item<1-> Hvad er sandsynligheden for at et positivt heltal valgt tilfældligt fra et sæt af positive heltal der ikke overgår 100 er deleligt med 2 eller 5?
  \item<1-> Hint: Definér $E_{1}$ og $E_{2}$ som værende tal der er delelige med hhv. 2 og 5.
  \item<2-> $E_{1}$ = tal der er delelige med 2, $E_{2}$ = tal der er delelige med 5.
  \item<3-> $|E_{1}| = 50$, $|E_{2}| = 20$, $|E_{1} \cup E_{2}| = 10$.
  \item<3-> Vi ved at $p(E_{1}\cup E_{2})=p(E_{1})+P(E_{2})-p(E_{1} \cap E_{2})$. Så: $= \frac{50}{100}+ \frac{20}{100}- \frac{10}{100} = \frac{3}{5}$
    
  \end{itemize} 
\end{frame}

\subsection{Probabilistic Reasoning}
\label{subsec:probabilistic-reasoning}

\begin{frame}
  \frametitle{Probabilistic Reasonining}
  \begin{itemize}
  \item Probabilistic Reasoning er at bruge sandsynlighedsregning til at ræsonnere sig frem til resultater. Et eksempel på hvor man kan gøre dette, er Monty Hall Three-Door Puzzle, som vi kommer til nu!
  \end{itemize} 
\end{frame}

\begin{frame}
  \frametitle{Eksempel 10: The Monty Hall Three-Door Puzzle}
  \begin{itemize}
  \item Antag at du er en spiller i et game show.
  \item Du bliver bedt om at vælge en af tre døre.
  \item Den store pris er bag en af de tre døre, og de to andre døre leder til tab.
  \item Når du vælger en dør, vil værten lade dig enten beholde døren, eller vælge en ny. Du får ikke lov til at se hvad der er i døren hvis du vælger en ny.
  \item Vælger du en ny? 
  \end{itemize}
\end{frame}

\begin{frame}
  \frametitle{Eksempel 10 Løsning}
  
  \begin{itemize}
  \item Ja du gør!
  \item Sandsynligheden for at du har valgt den rigtige dør først er $\frac{1}{3}$
  \item Sandsynligheden for at du har valg \textbf{forkert} er $\frac{2}{3}$. 
  \item Hvis du valgte forkert, åbner game show værten en dør hvor der ingen præmie er bag.
  \item Dermed, kommer du \textbf{altid} til at vinde, hvis dit første valg var forkert, og du så vælger at skifte døre.
  \item Så, ved at skifte døre er chancen for at vinde $\frac{2}{3}$
  \end{itemize}
\end{frame}

\begin{frame}
  \frametitle{Reflektion}
  \begin{itemize}
  \item Hvad er et eksperiment?
  \item Sample space?
  \item Event?
  \item Laplace's Definition?
  \item Hvad er værdimængden for $P(E)$? (Givet at $E$ er en event)
  \item Hvad er sampling with og without replacement?
  \item Hvordan fungerede beviset til $P(\overline{E})$?
  \end{itemize} 
\end{frame}

\section{Probability Theory}
\label{sec:probability-theortheory}

\subsection{Assigning Probabilities}
\label{subsec:assigning-probabilities}


\begin{frame}
  \frametitle{Assigning probabilities}
  We giver sandsynligheden $p(s)$ til hvert udfald $s$. We kræver at to betingelser bliver overholdt:
  \begin{itemize}
  \item $0 \leq p(s) \leq 1 \text{for each} s \in S$ og
  \item $\sum_{s \in S}p(s) = 1$
  \end{itemize}
  Funktionen $p$ bliver kaldt \textbf{probability distribution}.
\end{frame}

\begin{frame}
  \frametitle{Definition 1}
  \begin{definition}
    Suppose that $S$ is a set with $n$ elements. The \textit{uniform distribution} assigns the probability $1/n$ to each element of $S$.
  \end{definition}
  \begin{itemize}
  \item Basically, uniform distribution er hvor hvert udfald har lige stor sandsynlighed for at ske (og det er $1/n$ fordi vi ved at $\sum_{s \in S}^{}p(s) = 1$.)
  \end{itemize} 
\end{frame}

\begin{frame}
  \frametitle{Definition 2}
  \begin{definition}
    The \textit{probability}  of the event $E$ is the sum of the probabilities of the outcomes in E. That is $$p(E) = \sum_{s \in E}^{}p(s)$$

  \end{definition}
  \begin{itemize}
  \item Bare det jeg sagde på sidste slide, lol. Jeg er ikke helt sikker på hvorfor det er gentaget her.
  \end{itemize}
\end{frame}

\subsection{Probabilities of Complements and Unions of Events}
\label{subsec:label}

\begin{frame}
  \frametitle{Theorem 1}
  \begin{theorem}

    If $E_{1}, E_{2}, \ldots$ is a sequence of pairwise disjoint events in a sample space $S$, then $$p \left ( \bigcup_{i} E_{i} \right ) = \sum_{i}^{} p(E_{i})$$
  \end{theorem}
  \begin{itemize}
  \item Beviset undlades, men jeg mindes at det bringes op i KT 13.1
  \end{itemize}
\end{frame}

\subsection{Conditional Probability}
\label{subsec:condprob}

\begin{frame}
  \frametitle{Definition 3}
  \begin{definition}
    Let $E$ and $F$ be events with $p(F) > 0$. The \textit{conditional probability} of $E$ given $F$, denoted by $p(E |F)$, is defined as $$p(E|F) = \frac{p(E \cap F)}{p(F)}$$
  \end{definition}
\end{frame}

\begin{frame}
  \frametitle{Eksempel 3}

  \begin{itemize}
  \item<1-> En bit streng af længde 4 er tilfældligt generet, således at hver af de 16 bits af længde 4 er lige sandsynlige. Hvad er sandsynligheden at den har mindst to 0er i træk (Til oplægsholder: næste slide har lidt hjælp)
  \item<2-> Lad $E$ være begivenheden at en bit streng af længde 4 har mindst to 0er i træk. 
  \item<2-> Lad $F$ være begivenheden at den første bit af en streng af længde fire er et 0. 
  \item<2-> Hvad er sandsynligheden så?
  \item<3-> $p(E|F) = \frac{p(E \cap F)}{p(F)}$
  \item<3-> Siden $E \cap F = \{0000,0001,0010,0011,0100\}$, og dermed, $|E \cap F| = 5$, så må $p(E \cap F) = 5/16$. Ydermere er $p(F) = 1/2 = 8/16$. Derfor:
  \item<3-> $p(E|F) = \frac{5/16}{1/2}= \frac{5}{8}$
  \end{itemize}
\end{frame}

\subsection{Independence}
\label{subsec:independence}

\begin{frame}
  \frametitle{Independence}
  \begin{columns}
    \begin{column}{.5\textwidth}
      \begin{itemize}
      \item En begivenhed (event) er independent (uafhængig), hvis tidligere udfald ikke har effekt på sandsynligheden af begivenheden.
      \end{itemize}
      \begin{definition}[Definition 4]
        The events \(E\) and \(F\) are \textit{independent} if and only if \(p(E \cap F) = p(E)p(F)\)
      \end{definition}
    \end{column}
    \begin{column}{.5\textwidth}
      \includegraphics[width=\linewidth]{independence-day.png}
    \end{column}
  \end{columns} 
\end{frame}

\begin{frame}
  \frametitle{Eksempel 5}
  \begin{itemize}
  \item<1-> Antag at $E$ er begivenheden at en tilfældigt genereret big streng af længde 4 begynder med et 1, og $F$ er begivenheden at denne bit streng har præcis et lige antal 1'er. Er $E$ og $F$ uafhængige, hvis de 16 bits er lige sandsynlige?
  \item<2-> $E = \{100, 1001, 1010, 1011, 1100, 1101, 1110, 1111\}$
  \item<2-> $F = \{0000, 0011, 0101, 0110, 1001, 1010, 1100, 1111\}$.
  \item<3-> Dermed: $p(E) = p(F) = 8/16 = 1/2$
  \item<3-> Fordi $E \cap F = \{1111, 1100, 1010, 1001\}$, ser vi at:
  \item<3-> $p(E \cap F) = 4/16 = 1/4 = (1/2)(1/2) = p(E)p(F)$
  \item<3-> Dermed er $E$ og $F$ uafhængige.
  \end{itemize}
\end{frame}

\begin{frame}
  \frametitle{Pairwise and Mutual Independence}
  \begin{definition}[Definition 5]
    The events $E_{1}, E_{2}, \ldots, E_{n}$ are \textit{pairwise independent} if and only if $p(E_{i} \cap E_{j})=p(E_{i})p(E_{j})$ for all pairs of integers $i$ and $j$ with $1 \leq i< j \leq n$. These events are \textit{mutually independent} if $p(E_{i_{1}} \cap E_{i_{1}} \cap \cdots \cap E_{i_{m}}) = p(E_{i_{1}})p(E_{i_{2}}) \cdots p(E_{i_{m}})$ whenever $i_{j}, j = 1, 2, \ldots, m$ are integers with $1 \leq i_{1} < i_{2} < \cdots < i_{m} \leq n$ and $m \geq 2$.
  \end{definition}
  \begin{itemize}
  \item Forskellen er, i sin essens, at pairwise independence tillader at en begivenhed kan påvirke mere end én begivenheds sandsynlighed, bare ikke kun én.
  \item Derudover, kan sandsynligheder godt være mutually independent uden at være pairwise independent.
  \end{itemize}
\end{frame}

\subsection{Bernoulli Trials and the Binomial Distribution}
\label{subsec:bernoulli}

\begin{frame}
  \frametitle{Bernoulli Trials}
  \begin{itemize}
  \item En \textbf{Bernoulli Trial} er et ekseperiment hvor der kun er to mulige udfald. 
  \item Generelt er disse udfald kaldet \textit{success} eller \textit{fejl} (failure). 
  \item Givet at $p$ er sandsynligheden for success, og $q$ for fejl, må $p + q = 1$.
  \item Bernoulli Trials er mutually independent, hvis conditional probability af success på et givet trial er $p$, givet hvilket som helst information om udfaldet af andre \textit{trials}.
  \end{itemize} 
\end{frame}

\begin{frame}
  \frametitle{Eksempel 8}
  \begin{itemize}
  \item<1-> En mønt er biased, så sandsynligheden for plat er $2/3$. Hvad er sandsynligheden for at præcis fire plat kommer up når man kaster mønten syv gange? 
  \item<2-> Der er $2^{7} = 128$  forskellige udfald. Antallet af måder fire ud af de 7 møntkast kan være plat på er $\binom{7}{4}$. 
  \item<2-> Sandsynligheden for det præcise udfald er $(2/3)^{4}(1/3)^{3}$. Dermed er sandsynligheden:
  \item<3-> \[ \binom{7}{4} (2/3)^{4} (1/3)^{3} = \frac{35 \cdot 16}{3^{7}} = \frac{560}{2187} \]
  \end{itemize}
\end{frame}

\begin{frame}
  \frametitle{Theorem 2}
  \begin{theorem}[Theorem 2]
    The probability of exactly $k$ successes in $n$ independent Bernoulli trials, with probability of success $p$ and probability of failure $q = 1 - p$, is
    \[\binom{n}{k} p^{k} q^{n-k} \]
  \end{theorem}
\end{frame}
\begin{frame}
  \frametitle{Theorem 2 Bevis}
  Længe siden vi har set sådan et, hva? 
  \begin{proof}[Theorem 2 Bevis]
    Når $n$ Bernoulli trial bliver udført, er udfaldet en $n$-tuple $(t_{1}, t_{2},\ldots, t_{n})$, hvor $t_{i} = S$ (success) eller $t_{i} = F$ (fejl), for $i = 1, 2, \ldots, n$. Fordi de $n$ trials er uafhængige, så er sandsynligheden af hvert udfald af de $n$ trials af $k$ succeser og $n-k$ fejl $p^{k}q^{n-k}$. Fordi der er $\binom{n}{k}$ $n-$tupler af $S$'er og $F$'er som har præcis $k$ S'er, er sandsynligheden for præcis $k$ successer:
    \[ \binom{n}{k}p^{k}q^{n-k} \]
  \end{proof} 
\end{frame}

\begin{frame}
  \frametitle{Binomial Distribution}
  \begin{itemize}
  \item Vi betegner $b(k;n,p)$ til at være sandsynligheden af $k$ sucesser i $n$ uafhængige Bernoulli trials med sandsynlighed for succes $p$ og sandsynlighed for fejl $q = 1-p$. Vi kalder denne funktion \textbf{binomial distribution}. 
  \end{itemize}
\end{frame}

\subsection{Random Variables}
\label{subsec:random-variables}

\begin{frame}
  \frametitle{Random Variables}
  \begin{itemize}
  \item Random Variables, kendt som at være et af de mest vildledende navne i matematik efter Dynamic Programmering, er ellers ret vigtigt i sandsynlighedsregning.
  \item De er dog ret simple, de kommer nok ikke til at fylde en specielt stor del af resten af kurset.
  \end{itemize} 
\end{frame}

\begin{frame}
  \frametitle{Definition 6}
  \begin{definition}
    A \textit{random variable} is a function from the sample space of an experiment to the set of real numbers. That is, a random variable assigns a real number to each posssible outcome.
  \end{definition}
  \begin{itemize}
  \item Så.. Det er bare en tilfældig variabel?
  \item Nej!
  \item Den er hverken tilfældig, og den er heller ikke variabel! Det er en funktion!!!
  \end{itemize}
\end{frame}

\begin{frame}
  \frametitle{Definition 7}
  \begin{definition}
    The \textit{distribution} of a random variable $X$ on a sample space $S$ is the set of pairs $(r, p(X=r))$  for all $r \in X(S)$, where $p(X=r)$ is the probability that $X$ takes the value $r$. (The set of pairs in this distribution is determined by the probabilities $p(X=r)$ for $r \in X(S)$.)
  \end{definition}
\end{frame}

\begin{frame}[allowframebreaks]
  \frametitle{Eksempel 13: The Birthday Problem}
  \begin{itemize}
  \item (Dette er ikke en opgave lignende andre eksempler, jeg giver svaret, og forventer ikke at I kan det).
  \item Hvad er det minimum antal af folk der skal være i et rum således at sandsynligheden for at to af dem har samme fødselsdag er større end $1/2$?
  \item Vi antager at fødselsdagene er uafhængige. 
  \item Vi antager at hver fødselsdag har lige sandsynlighed for at være præcis den dag, i.e. sandsynligheden for at fødselsdagen er på en præcis dag: $1/n$ hvor $n$ er antallet af dage, som er:
  \item Vi antager at antallet af dage på et år er 366.
  \item Fremgangsmåde:
    \begin{itemize}
    \item Vi må først udregne sandsynligheden $p_{n}$ for at all $n$ personer har forskellige fødselsdage
    \item Derefter taget vi $1-p_{n}$, som så giver os sandsynligheden for at mindst to personer deler fødselsdag.
    \end{itemize}
  \item Sandsynligheden for at to personer har samme fødselsdag er $\frac{365}{366}$. Dermed, ved at vi tager til $j$ personer bliver det:
  \item \[ \frac{366(j-1)}{366} = \frac{367-j}{366} \]
  \item Fordi vi har antaget at fødselsdagene af personerne i rummet er uafhængige, kan vi konkludere at sandsynligheden at de $n$ personer i rummet har forskellige fødselsdage er: \[ p_{n} = \frac{365}{366} \frac{364}{366} \frac{363}{366} \cdots \frac{367-n}{366}\].
  \item Sandsynligheden for at to har samme fødselsdag er $1-p_{n}$.
  \item Vi gider ikke til at bruge differentialregning eller whatever, så vi prøver os bare frem, og finder ud af at denne formel giver $n = 23, 1-p_{n} \approx 0.506$.
  \end{itemize}
\end{frame}

\begin{frame}[allowframebreaks]
  \frametitle{Eksempel 14: Probability of a Collision in Hashing Functions}
  \begin{itemize}
  \item Denne her er god at kunne når vi kommer til hashing!
  \item Hvad er sandsynligheden for at ingen nøgler mapper til den samme lokation af én hashing funktion?
  \item Vi antager universal hashing function (i.e., sandsynligheden for at den rammer en specifik lokation er $1/m$, hvor $m$ er antallet af slots.)
  \item Vi modellerer nøglerne til at være $k_{1}, k_{2}, \ldots, k_{n}$, hvor $n$ er antallet af nøgler.
  \item Sandsynligheden for at $h(k_{2}) \neq h(k_{1}) = \frac{(m-1)}{m}$ fordi der er $m-1$ lokationer den kan komme til, ud af $m$ i alt, og den første er allerede taget. Det fortsætter og bliver til:
  \item $\frac{m-(j-1)}{m}$ hvor $j$ er nøglenummeret (fra 1 opad). Dermed, givet at $p_{n}$ er sandsynligheden for at $h(k_{1}) \neq h(k_{2}) \neq \cdots \neq h(k_{n}) = p_{n}$, \[p_{n}= \frac{m-1}{m} \cdot \frac{m-2}{m } \cdot \cdots \cdot \frac{m-n+1}{m} \]. 
  \item Vi finder så sandsynligheden ved at tage komplement, i.e. $1 - p_{n}$.
  \end{itemize}
\end{frame}

\subsection{The Probabilistic Method}
\label{subsec:probabilistic-method}

\begin{frame}
  \frametitle{Theorem 3}
  \begin{theorem}[Theorem 3]
    \textbf{The Probabilistic Method} If the probability that an element chosen at random from a $S$ does not have a particular property is less than 1, there exists an elementi n $S$ with this property.
  \end{theorem}
  \begin{itemize}
  \item Kan blive brugt til at lave nonkonstruktive eksistensbeviser (i.e., du viser at et element eksisterer, du viser bare ikke hvor det er, eller hvordan du finder det).
  \end{itemize}
\end{frame}

\begin{frame}
  \frametitle{Theorem 4}
  \begin{theorem}
    If $k$ is an integer with $k \geq 2$, then $R(k,k) \geq 2^{k/2}.$
  \end{theorem}
\end{frame}

\begin{frame}[allowframebreaks]
  \frametitle{Theorem 4 Bevis}
  Vi ved at teoremet holder for $k = 2$ og $k = 3$, fordi $R(2,2) = 2, R(3,3) = 6$. (Dette var vist i kapitel 6.2).

  Hvis vi så antager at $k \geq 4$, vil vi bruge den probabilistike metode til at vise om der er færre end $2^{k/2}$ personer til en fest, så er det muligt at ingen $k$ af dem er \textit{mutual friends or mutual enemies} (jeg gider ikke til at oversætte).

  Vi antager at sandsynligheden for at to personer er venner eller fjender er lige stor.

  Antag at der er $n$ personer til festen. Dermed er der $\binom{n}{k}$ forskellige sæt af $k$ personer til denne færst, hvilke vi skriver som $S_{1}, S_{2}, \ldots, S_{\binom{n}{k}}$ (altså, sættene af personer).

  Lad $E_{i}$ være begivenheden at alle $k$  personer i $S_{i}$ er enten \textit{mutual friends} eller \textit{mutual enemies}.

  Sandsynligheden for at der enten er $k$ \textit{mutual friends} eller \textit{mutual enemies} blandt de $n$ personer er lig med \[ \bigcup\limits_{i=1}^{\binom{n}{k}}E_{i} \].

  Ifølge vores antagelse at det er lige sandsynligt for to personer at være fjender eller venner, så er sandsynligheden for begge $1/2$. Ydermere er der $\binom{k}{2} = \frac{k(k-1)}{2}$ par af personer i $S_{i}$, fordi der er $k$ personer i $S_{i}$. Dermed er sandsynligheden for at alle $k$ personer i $S_{i}$ er \textit{mutual friends} og sandsynligheden for at alle k personer i $S_{i}$ er \textit{mutual enemies} er begge lig $(1/2)^{k(k-1)/2}$. Det følger at $p(E_{i}) = 2(1/2)^{k(k-1)/2}$.

  Sandsynligheden for at der er enten $k$ \textit{mutual friends} eller $k$ \textit{mutual enemies} i gruppen af $n$ personer er lig med $p(\bigcup\limits_{i=1}^{\binom{n}{k}} E_{i})$. Ved brug af Boole's Inequality, følger det at:

  \[ p\left( \bigcup\limits_{i=1}^{\binom{n}{k}} \leq \sum_{i=1}^{\binom{n}{k}} p(E_{i}) = \binom{n}{k} \cdot 2 \left( \frac{1}{2} \right)^{k(k-1)/2}   \right) \]

  Af exercise 21 i 6.4 har vi at $\binom{n}{k} \leq n^{k}/2^{k-1}$. Dermed,
  \[ \binom{n}{k} 2 \left( \frac{1}{2} \right)^{k(k-1)/2} \leq \frac{n^{k}}{2^{k-1}}2 \left( \frac{1}{2} \right)^{k(k-1)/2} \]

  Hvis $n < 2^{k/2}$, har vi at :

  \[ \frac{n^{k}}{2^{k-1}}2 \left( \frac{1}{2} \right)^{k(k-1)/2} < \frac{2^{k(k/2)}}{2^{k-1}}2 \left( \frac{1}{2} \right)^{k(k-1)/2} = 2^{2-(k/2} \leq 1 \]
\end{frame}

\begin{frame}[allowframebreaks]
  \frametitle{Reflektion}
  \begin{itemize}
  \item Hvad er probability distribution? 
  \item Hvad er en uniform distribution?
  \item Hvad er conditional probability?
  \item Hvordan udregnes $P(E|F)$?
  \item Forklar, uden brug af matematik, hvad uafhængige begivenheder er.
  \item Forklar med matematik.
  \item Hvad er pairwise independence
  \item Hvad er mutual independence
  \item Hvad er forskellen?
  \item Ligheden?
  \item Hvad er en bernoulli trial?
  \item Hvad er en binomial distribution (binomial fordeling)?
  \item Hvad er en random variable? 
  \item Er det en variable?
  \item Hvad er den probabilistiske metode?
  \end{itemize} 
\end{frame}

\section{Bayes' Theorem}
\label{sec:bayes}

\begin{frame}
  \frametitle{Bayes' Theorem}
  \begin{itemize}
  \item Bayes' Theorem er et teorem der kan hjælpe dig med at finde en sandsynlighed, baseret udelukkende på partial knowledge.
  \end{itemize} 
\end{frame}

\subsection{Bayes' Theorem}
\label{subsec:bayes-theorem}


\begin{frame}[allowframebreaks]
  \frametitle{Eksempel 1}

  \begin{itemize}
  \item Vi har to bokse. Den første indeholder to grønne bolde, og syv røde bolde. Den anden indeholder fire grønne bolde og tre grønne bolde.
  \item Bob vælger en bold ved først at vælge en af de to bokse tilfældigt, og derefter en bold tilfældeligt. Hvis bob har valgt en rød bold, hvad er sandsynligheden at han tog fra den første boks? 
  \item Lad $E$ være begivenheden at Bob har valgt en rød bold. $\overline{E}$ er begivenheden at Bob har valgt en grøn bold. Lad $F$ være begivenheden at Bob har valgt fra den første boks ($\overline{F}$ er dermed begivenheden at han har valgt fra anden).
  \item Vi vil gerne finde $P(F|E)$. Vi ved at $P(F|E) = \frac{P(F \cap E)}{P(E)}$. Kan vi bruge informationen givet til at finde både $P(F \cap E)$ og $P(E)$, så vi kan finde $P(F|E)$?
  \item For det første ved vi at vi kan finde sandsynligheden for at han vælger en rød bold, hvis han tager fra den første boks. $P(E|F) = 7/9$. Vi ved også sandsynligheden for at han tager en rød bold hvis han tager fra den anden boks: $P(E|\overline{F}) = 3/7$. 
  \item Siden vi antager at Bob vælger en boks tilfældig, er $P(F) = P(\overline{F}) = 1/2$.
  \item Vi ved at $p(E|F) = \frac{p(E \cap F)}{p(F)}$. Hvis vi ganger med $p(F)$ på begge sider får vi: $p(E|F) \cdot p(F) = \frac{p(E \cap F)}{p(F)} \cdot p(F) = p(E \cap F)$. Dermed kan vi finde $p(E \cap F)$ ved at have $p(E|F)$ og $p(F)$. Vi har etableret, at vi har denne information.
  \item $\frac{7}{9}\cdot \frac{1}{2} = \frac{7}{18}$
  \item Ligeledes: $p(E \cap \overline{F}) = p(E| \overline{F}) p(\overline{F}) = \frac{3}{7} \cdot \frac{1}{2} = \frac{3}{14}$
  \item Da $E = (E \cap F) \cup (E \cap \overline{F})$ kan vi nu finde $p(E)$.
  \item $p(E) = p(E \cap F) + p(E \cap \overline{F}) = \frac{7}{18} + \frac{3}{14} = \frac{49}{126} + \frac{27}{126} = \frac{76}{126} = \frac{38}{63}$
  \item Dermed:
  \item \[ p(F|E) = \frac{p(F \cap E)}{p(E)} = \frac{7/18}{38/63} = \frac{49}{76} \approx 0.645 \]
  \end{itemize}
\end{frame}

\begin{frame}
  \frametitle{Theorem 1}
  \begin{theorem}[Bayes' Theorem]
    Suppose that $E$ and $F$ are events from a sample space $S$ such that $p(E) \neq 0$ and $p(F) \neq 0$. Then
    \[ p(F|E) = \frac{p(E|F)p(F)}{p(E|F)p(F)+p(E|\overline{F})p(\overline{F})} \]
  \end{theorem} 
\end{frame}

\begin{frame}[allowframebreaks]
  \frametitle{Bayes' Theorem Bevis}
  Definitionen af conditional probability fortæller os at $p(F|E) = p(E \cap F)/p(E)$ og $p(E|F) = p(E \cap F)/p(F)$. Derfor, $p(E \cap F) = p(F|E)p(E)$ og $p(E \cap F) = p(E|F)p(F)$. Hvis vi sætter disse lig hinanden, får vi:
  \[ p(F|E)p(E) = p(E|F)p(F)  \]

  Hvis vi dividerer begge disse sider med $p(E)$ får vi:
  \[ p(F|E) = \frac{p(E|F)p(F)}{p(E)} \]

  Som det næste viser vi at $p(E) = p(E|F)p(F)+P(E|\overline{F})p(\overline{F})$. For at se dette, så læg først mærke til at $E = E \cap S = E \cap (F \cup \overline{F}) = (E \cap F) \cup ( E \cap \overline{F})$. Ydermere, $E \cap F$ og $E \cap \overline{F}$ er \textit{disjoint} (usammenhængende), fordi hvis $x \in E \cap F$ og $x \in E \cap \overline{F}$, så $x \in F \cap \overline{F} = \emptyset$.

  Konsekvent må $p(E) = p(E \cap F) + p(E \cap \overline{F})$. Vi har allerede vist at $p(E \cap F) = p(E |F) p(F)$. Derudover, har vi at $p(E | \overline{F}) = p(E \cap \overline{F})/p(\overline{F})$, som viser os at $p(E \cap \overline{F}) = p(E|\overline{F})p(\overline{F})$. Det følger nu at:
  \[ p(E) = p(E \cap F) + p(E \cap \overline{F}) = p(E |F) p(F) + p(E | \overline{F} p(\overline{F}) \]

  For at færdiggøre beviset plopper vi bare vores nye $p(E)$ ind og får:
  \[ p(F|E) = \frac{p(E|F)p(F)}{p(E|F)p(F)+p(E|\overline{F})p(\overline{F})} \]
\end{frame}

\begin{frame}
  \frametitle{Theorem 2}
  \begin{theorem}[Generalized Bayes' Theorem]
    Suppose that $E$ is an event from a sample space $S$ and that $F_{1}$, $F_{2}, \ldots, F_{n}$ are mutually exclusive events such that $\bigcup\limits_{i=1}^{n}F_{i}=S$. Assume that $p(E) \neq 0$ and $p(F_{i}) \neq 0$, for $i = 1, 2, \ldots, n$. Then
    \[ p(F_{j}|E) = \frac{p(E|F_{j}) p(F_{j})}{\sum_{i=1}^{n}p(E|F_{i})p(F_{i})} \]
  \end{theorem}
\end{frame}

\begin{frame}
  \frametitle{Reflektion}
  \begin{itemize}
  \item Hvad er Bayes' Theorem?
  \item Hvad kan det bruges til?
  \item Hvornår?
  \item Hvorfor er $E = E \cap (F \cup \overline{F})$?
  \end{itemize}
  
\end{frame}

\subsection{Expected Value and Variance}
\label{subsec:expectedvalue}

\begin{frame}
  \frametitle{Expected Value}
  \begin{itemize}
  \item \textbf{Expected Value} af en random variable er summen over alle elementer i et sample space af produktet af sandsynligheden af elementet og værdien af den tilfældelige variable på dette element. Konsekvent er expected value et vægtet gennemsnit af værdierne af en random variable.
  \end{itemize}
\end{frame}

\subsection{Expected Values}
\label{subsec:expected-values}

\begin{frame}
  \frametitle{Definition 1}
  \begin{definition}
    The \textit{expected value}, also called the \textit{expectation} or \textit{mean}, of the random variable $X$ on the sample space $S$ is equal to
    \[ E(X) = \sum_{s \in S}^{}p(s)X(s) \]
    The \textit{deviation} of $X$ at $s \in S$ is $X(s)-E(X)$, the difference between the value of $X$ and the mean of $X$.
  \end{definition}
  Læg mærke til at når sample space $S$ har $n$ elementer, $S = \{x_{1}, x_{2}, \ldots, x_{n}\}, E(X) = \sum_{i=1}^{n}p(x_{i})X(x_{i})$
\end{frame}

\begin{frame}
  \frametitle{Eksempel 1}
  \begin{itemize}
  \item<1-> Nu er vi tilbage til nogle lidt nemmere opgaver. 
  \item<1-> Så her er der også nogle skjulte svar.
  \item<1-> Lad $X$ være tallet der kommer up når en terning bliver rullet. Hvad er expected value af $X$? (Næste slide indeholder sandsynligheden, ikke $E(X)$)
  \item<2-> $p(X = r) = 1/6, r = 1,2,\ldots 6$
  \item<3-> $E[X] = \frac{1}{6} \cdot 1 + \frac{1}{6} \cdot 2 + \cdots + \frac{1}{6} \cdot 6 = \frac{21}{6} = \frac{7}{2}$
  \end{itemize} 
\end{frame}


\begin{frame}
  \frametitle{Theorem 1}
  \begin{theorem}
    If $X$ is a random variable and $p(X = r)$ is the probability that $X = r$, so that $p(X=r) = \sum_{s \in S, X(s)=r} p(s)$, then
    \[ E[X] = \sum_{r \in X(S)}^{} p(X=r)r  \]
  \end{theorem} 
  \begin{itemize}
  \item a.k.a., den langt nemmere måde at udregne $E[X]$ på.
  \item vent? er der en forskel? er oprigtigt i tvivl
  \end{itemize}
\end{frame}

\begin{frame}
  \frametitle{Theorem 1 Bevis}
  \begin{proof}
    (Endelig et bevis så kort at vi kan have bevis-rammen tilbage!)
    Antag at $X$ er en random variable med værdimængde $X(S)$, og let $p(X=r)$ være sandsynligheden at en random variable $X$ tager vrædien $r$. Konsekvent, $p(X=r)$ er summen af sandsynlighederne $s$, således at $X(s) = r$. Det følger at
    \[E[X] = \sum_{r \in X(S)}^{} p(X = r)r  \]
  \end{proof}
\end{frame}

\begin{frame}
  \frametitle{Forskel}
  \begin{itemize}
  \item Ifølge ChatGPT er forskellene at:
    \begin{itemize}
    \item Definition 1 summerer over hele sample spacet, og udregner indsatsen af hvert udfald individuelt.
    \item Teoremet summerer over de distinkte værdier som et random variable kan have. 
    \end{itemize}
  \end{itemize}
\end{frame}


\begin{frame}
  \frametitle{Theorem 2}
  \begin{theorem}
    The expected number of successes when $n$ mutually independent \textbf{Bernoulli Trials} are performed, where $p$ is the probability of success on each trial, is $np$.
  \end{theorem} 
\end{frame}

\begin{frame}
  \frametitle{Theorem 2 Bevis}
  Lad $X$  være det random variable lig med antallet af sucesser i $n$ trials. Af Teorem 2 i 7.2 ser vi at $p(X=k) = \binom{n}{k}p^{k}q^{n-k}$. Dermed har vi:

\end{frame}

\begin{frame}
  \frametitle{Theorem 2 bevis part 2}
  \begin{equation*}
    \begin{split}
      E(X) &= \sum_{k=1}^{n}kp(X=k) \;\; \text{by Theorem 1}\\
           &= \sum_{k=1}^{n}k \binom{n}{k}p^{k}q^{n-k} \;\; \text{by Theorem 2 in 7.2}\\
           &= \sum_{k=1}^{n} n \binom{n-1}{k-1}p^{k} q^{n-k} \; \; \text{by Exercise 21 in 6.4}\\
           &= np\sum_{k=1}^{n} \binom{n-1}{k-1}p^{k-1}q^{n-k} \; \; \text{factoring np from each term} \\
           &= np \sum_{j=0}^{n-1} \binom{n-1}{j}p^{j}q^{n-1-j}\;\; \text{shifting index of summation with} j = k - 1\\
           &= np(p+q)^{n-1} \;\; \text{by the binomial theorem}\\
           &= np \;\; \text{because p + q = 1}\\
    \end{split}
  \end{equation*}
\end{frame}

\subsection{Linearity of Expectations}
\label{subsec:linexp}

\begin{frame}
  \frametitle{Theorem 3}
  \begin{theorem}
    If $X_i, i = 1, 2, \ldots, n$ with $n$ a positive integer, are random variables on $S$, and if $a$ and $b$ are real numbers, then
    \begin{itemize}
    \item $E(X_{1}+X_{2}+ \cdots + X_{n}) = E(X_{1})+ E(X_{2})+ \cdots E(X_{n})$
    \item $E(aX+b) = aE(X)+b$
    \end{itemize}
  \end{theorem}
\end{frame}


\begin{frame}
  \frametitle{Theorem 3 Beviser}
  Part 1:
  \begin{equation}
    \begin{split}
      E(X_{1} + X_{2}) &= \sum_{s \in S}p(s) (X_{1}(s) + X_2(s)) \\
                       &= \sum_{s \in S}p(s) X_{1}(s) + \sum_{s \in S}p(s) X_{2}(s)\\
                       &= E(X_{1}) + E(X_{2})
    \end{split}
  \end{equation}

  Part 2:
  \begin{equation}
    \begin{split}
      E(aX+b) &= \sum_{s \in S}p(s)(aX(s)+b)\\
              &= a \sum_{s \in S}p(s)X(s) + b \sum_{s \in S}p(s)\\
              &= aE(X) + b \text{ because } \sum_{s \in S}p(s) = 1\\
    \end{split}
  \end{equation}
\end{frame}

\begin{frame}
  \frametitle{Eksempel 4}
  \begin{itemize}
  \item<1-> Brug teorem 3 til at finde expected value af summen af tallene der fremkommer når et par af terninger bliver ullet. 
  \item<1-> Hint: Lad $X_{1}$  og $X_{2}$ være random variables, hvor $X_{1}((i,j)) = i$, og $X_{2}((i,j)) = j$, altså, $X_{1}$ er værdien på første terning, og $X_{2}$ på anden terning. Hvad er værdierne af $X_{1}$ og $X_{2}$?
  \item<2-> Herfra kan det ses at $E[X_{1}] = E[X_{2}] = 7/2 = (1+2+3+4+5+6)/6$.
  \item<3-> Ved teorem 3: $E[X_{1}+X_{2}] = E[X_{1}] + E[X_{2}] = 7/2 + 7/2 = 7$
  \end{itemize} 
\end{frame}

\begin{frame}
  \frametitle{Expected Value In The Hatcheck Problem}
  \begin{itemize}
  \item<1-> Husker I Hatcheck Problemet? Lad os finde expected value.
  \item<1-> Hint: Lad $X$ være random variable der er lig med antallet af folk som får den korrekte hat fra tjekkeren. Lad $X_{i}$ være random variable hvor $X_{i} = 1$ hvis den i'e person får den korrekte hat, og $X_{i}=0$ ellers (independent variables foreshadowing much?)
  \item<1-> Dermed må $X = X_1+ X_{2}+ \ldots+ X_{n}$. Næste: udregn sandsynlighed for $P(X_{i}=1)$
  \item<2-> $P(X_{i}=1) = 1/n$. Hvad er så $E[X_{i}]$?
  \item<3-> $E[X_{i}] = 1 \cdot p(X_{i}=1) + 0 \cdot p(X_{i}=0) = 1/n$. Hvad er så $E[X]$?
  \item<4-> $E[X] = E[X_{1}] + E[X_{2}]+ \cdots + E[X_{n}] = n \cdot 1/n = 1$
  \end{itemize}
\end{frame}


\begin{frame}[allowframebreaks]
  \frametitle{Eksempel 7: Expected Number of Inversions in a Permutation}
  \begin{itemize}
  \item The ordered pair $(i,j)$ is called an \textit{inversion} in a permutation of the first $n$ positive integers if $i < j$, but $j$ precedes $i$ in the permutation. For instance, there are six inversions in the permutation $3,5,1,4,2$. These are:
  \item $(1,3),(1,5),(2,3),(2,4),(2,5),(4,5)$
  \item Let $I_{i,j}$ be the random variable on the set of all permutations of the first $n$ positive integers with $I_{i,j} = 1$ if $(i,j)$ is an inversion of the permutation and $I_{i,j} = 0$ otherwise. It follows that if $X$ is the random variable equal to the number of inversions in the permutation, then \[\sum_{1 \leq i < j \leq n}^{} I_{i,j}  \]
  \item Note that it is equally likely for $i$ to precede $j$ in a randomly chosen permutation as it is for $j$ to precede $i$. (To see this, note that there are an equal number
  \item Note that it is equally likely for $i$ to precede $j$ in a randomly chosen permutation, as it is for $j$ to precede $i$. Therefore:
  \item $E[I_{i,j}] = 1 \cdot p(I_{i,j}=1)+0 \cdot p(I_{i,j}=0)=1 \cdot 1/2 + 0 = 1/2$
  \item Because there are $\binom{n}{2}$ pairs $i$ and $j$ with $1 \leq i < j \leq n$ and by the linearity of expectations, we have \[ E[X] = \sum_{1 \leq i < j \leq n}^{} E[I_{i,j}] = \binom{n}{2} \cdot \frac{1}{2} = \frac{n(n-1)}{4} \]
  \end{itemize}
\end{frame}

\subsection{Average-Case Computational Complexity}
\label{subsec:average}

\begin{frame}
  \frametitle{Average-case Computational Complexity}
  \begin{itemize}
  \item Du kan se det som at udregne expected value af en random variable. 
  \item Lad sample space af et ekseperiment være sættet af mulige inputs, $a_{j}, j = 1, 2, \ldots, n$, og lad $X$ være det random variable der giver $a_{j}$ antallet af operationer brugt af algoritmen når $a_{j}$ er brugt som input. Vi giver sandsynligheden $p(a_{j})$ to hvert input værdi. Så er average-case complexity af algoritmen:
  \item $E[X] = \sum_{j=1}^{n} p(a_{j}) X(a_{j})$
  \end{itemize} 
\end{frame}


\subsection{The Geometric Distribution}
\label{subsec:geometric-distribution}

\begin{frame}
  \frametitle{Eksempel 10}
  \begin{itemize}
  \item Antag at sandsynligheden for at en mønt lander på plat er $p$. Mønten bliver kastet indtil den rammer plat. Hvor mange gange skal man kaste mønten før man får plat?
  \item Sandsynligheden for plat er $p(P) = p$. Videre: $p(KP) = (1-p)p, p(KKP) = (1-p)^{2}p$ etc.
  \item Generelt, med $n$ antal kroner, er sandsynligheden $(1-p)^{n-1}p$.
  \item \[ E[X] = \sum_{j=1}^{\infty}j\cdot p(X = j) = \sum_{j=1}^{\infty} j(1-p)^{j-1} p = p \sum_{j=1}^{\infty} j(1-p)^{j-1} = p \cdot \frac{1}{p^{2}} = \frac{1}{p} \]
  \end{itemize}
\end{frame}

\begin{frame}
  \frametitle{Geometric Distrbution}

  Et random variable hvor antallet af kast før en mønt bliver plat, er et random variable med \textbf{geometric distribution} 
  \begin{definition}
    A random variable $X$ has a \textit{geometric distribution with parameter } $p$ if $p(X=k)=(1-p)^{k-1}p$ for $k=1,2,3, \ldots,$ where $p$ is a real number with $0 \leq p \leq 1$.
  \end{definition}

  \begin{theorem}[Theorem 4]
    If the random variable $X$ has the geometric distribution with parameter $p$, then $E[X] = 1/p$
  \end{theorem}
\end{frame}

\subsection{Independent Random Variables}
\label{subsec:indeprandvar}

\begin{frame}
  \frametitle{Definition 3}

  \begin{definition}[Definition 3]
    The random variables $X$ and $Y$ on a sample space $S$ are \textit{independent} if
    \[ p(X=r_{1} \text{ and } Y = r_{2}) = p(X=r_{1}) \cdot p(Y = r_{2}) \]
    or in words, if the probability that $X=r_{1}$ and $Y=r_{2}$ equals the product of the probabilities that $X = r_{1}$ and $Y = R_{2}$, for all real numbers $r_{1}$ and $r_{2}$.
  \end{definition}
\end{frame}

\begin{frame}[allowframebreaks]
  \frametitle{Eksempel 11}
  \begin{itemize}
  \item Er random variables $X_{1}$ og $X_2$ fra eksempel 4 uafhængige (independent)?
  \item Eksempel 4 er eksemplet med terningekast.
  \item Lad $S = \{1,2,3,4,5,6\}$, og lad $i \in S$ og $j <in S$. Fordi der er $36$ forskellige udfald når parret af terninger er rullet, og hvert er lige sandsynligt, har vi:
    \[ p(X_{1}=i \text{ and } X_{2} =j) = 1/36 \].
  \item Ydermere, $p(X_{1}=i = 1/6$ og $p(X_{2}=j) = 1/6$, fordi sandsynligheden $i$ findes på den første terning, og sandsynligheden at $j$ findes på den anden terning er begge $1/6$, følger det at:
    \[ p(X_{1} = i) \text{ and } p(X_{2}=j) = \frac{1}{36} \]
    og
    \[p(X_{1} = i)p(X_{2}=j) = \frac{1}{6}\cdot \frac{1}{6} = \frac{1}{36} \]
  \item Dermed er $X_1$ og $X_{2}$ uafhængige.
  \end{itemize}
\end{frame}

\begin{frame}[allowframebreaks]
  \frametitle{Eksempel 12}
  \begin{itemize}
  \item Vis at random variablerne $X_1$ og $X = X_{1} + X_{2}$, hvor $X_{1}$ og $X_{2}$ er som defineret i eksempel 4, ikke er uafhængige. 
  \item Se først at $p(X_{1} = 1 \text{ og } X = 12 = 0$, fordi $X_{1}$ betyder at antallet der vises på den første terning er 1, hvilket indebærer at summet ikke kan være 12. 
  \item På den anden hånd er $p(X_{1}=1)=1/6$ og $p(X=12) = 1/36$. Dermed $p(X_{1}=1 \text{ og } X = 12) \neq p(X_{1} =1) \cdot p(X=12)$. Dermed er de ikke uafhængige.
  \end{itemize}
\end{frame}

\begin{frame}
  \frametitle{Theorem 5}
  \begin{theorem}[Theorem 5]
    If $X$ and $Y$ are independent random variables on a sample space $S$, then $E[XY] = E[X]E[Y]$.
  \end{theorem}
\end{frame}

\begin{frame}[allowframebreaks]
  \frametitle{Theorem 5 Bevis}
  For at bevise dette bruger vi nøgleobservationen at $XY = r$  er den disjunkte forening af hændelserne $X = r_{1}$ og $Y= r_{2}$ hvor $r_{1} \in X(S)$ og $r_{2} \in Y(S)$ med $r = r_{1}r_{2}$. Vi har

\end{frame}

\begin{frame}
  \frametitle{Theorem 5 Bevis part 2}
  \begin{equation*}
    \begin{split}
      E[XY] &= \sum_{r \in XY(S)}^{} r \cdot p(XY = r)\;\; \text{by Theorem 1}\\
            &= \sum_{r_{1} \in X(S), r_{2} \in Y(S)}^{} r_{1} r_{2} \cdot p(X = r_{1} \text{ og } Y = r_{2})\\
            &= \sum_{r_{1} \in X(S)}^{} \sum_{r_{2} \in Y(S)}^{} r_{1} r_{2} \cdot p(X = r_{1} \text{ og } Y = r_{2})\\
            &= \sum_{r_{1} \in X(S)}^{} \sum_{r_{2} \in Y(S)}^{} r_{1} r_{2} \cdot p(X = r_{1} \cdot Y = r_{2})\\
            &= \sum_{r_{1} \in X(S)}^{} \left(  r_{1} \cdot p(X=r_{1}) \cdot \sum_{r_{2} \in Y(S)}^{} r_{2} \cdot p(Y=r_{2}) \right)\\
    \end{split}
  \end{equation*}
\end{frame}


\begin{frame}
  \frametitle{Theorem 5 Proof Part 3}
  
  \begin{equation*}
    \begin{split}
      &= \sum_{r_{1} \in X(S)} r_{1} \cdot p(X = r_{1}) \cdot E(Y)\\
      &= E(Y) \left( \sum_{r_{1} \in X(S)} r_{1} \cdot p(X = r_{1}) \right)\\
      &= E(Y) \cdot E(X)
    \end{split}
  \end{equation*}
\end{frame}


\subsection{Variance}
\label{subsec:variance}

\begin{frame}
  \frametitle{Variance}
  Variansen er en variabel der bestemmer hvor meget en random variable spreder sig. Ikke at jeg forstår hvad det betyder, men altså.
  \begin{definition}
    Let $X$ be a random variable on a sample space $S$. The \textit{variance} of $X$, denoted by $V(X)$, is
    \[ V(X) = \sum_{s\in S}^{} (X(s) - E(X))^{2}p(s) \]
    That is, $V(X)$ is the weighted average of the square of the deviation of $X$. The \textit{standard deviation} of $X$, denoted $\delta(X)$, is defined to be $\sqrt{V(X)}$.
  \end{definition}
  
\end{frame}

\begin{frame}
  \frametitle{Theorem 6}
  \begin{theorem}[Theorem 6]
If $X$ is a random variable on a sample space $S$, then $V(X) = E(X^{2}) -E(X)^{2}$.
  \end{theorem}
\end{frame}

\begin{frame}
  \frametitle{Theorem 6 Bevis}
  Læg mærke til at:
  \begin{equation*}
    \begin{split}
      V(X) &= \sum_{s \in S}^{} (X(S) - E(X))^{2}p(s)\\
           &= \sum_{s \in S}^{} X(s)^{2}p(s) - 2E(X) \sum_{s \in S}^{} X(s)p(s) + E(X)^{2} \sum_{s \in S}^{} p(s)\\
           &= E(X^{2}) - 2E(X)E(X)+E(X)^{2}\\
           &= E(X^{2})-E(X)^{2}\\
    \end{split}
  \end{equation*}
\end{frame}


\begin{frame}
  \frametitle{Corollary 1}
  \begin{corollary}
    If $X$ is a random variable on a sample space $S$ and $E(X) = \mu$, then $V(X) = E((X- \mu)^{2})$.
  \end{corollary} 
\end{frame}

\begin{frame}
  \frametitle{Corollary 1 Proof}
  \begin{equation*}
    \begin{split}
      E((X-\mu)^{2}) &= E(X^{2}-2 \mu X + \mu ^{2})\\
                   &= E(X^{2})-E(2 \mu X) + E(\mu ^{2})\\
                   &= E(X^{2}) - 2 \mu E(X) + E(\mu ^{2})\\
                   &= E(X^{2}) - 2 \mu E(X) + \mu ^{2}\; \text{because } \mu ^{2} \text{ is a constant}\\
                   &= E(X^{2}) - 2 \mu^{2} + \mu^{2}\\
                   &= E(X^{2}) - \mu^{2}\\
                   &= V(X)\\
    \end{split}
  \end{equation*}
\end{frame}


\begin{frame}
  \frametitle{Eksempel 15: Variance of the Value of a Die}
  \begin{itemize}
  \item Hvad er variansen af random variable $X$, hvor $X$ er tallet der kommer op når en terning bliver kastet?
  \item Svar: $V(X) = E(X^{2})-E(X)^{2}$. Vi ved at $E(X) = 7/2$. $E(X^{2}) = \frac{1}{6}(1^{2}+2^{2} + \cdots + 6^{2})= \frac{91}{6}$. 
  \item Dermed: 
  \item $V(X) = \frac{91}{6} - \left( \frac{7}{2}\right)^{2} = \frac{35}{12}$
  \end{itemize}
\end{frame}


\begin{frame}
  \frametitle{Bienaymé's Formula}
  \begin{theorem}
    If $X$ and $Y$ are two independent random variables on a sample space $S$, then $V(X+Y) = V(X) + V(Y)$. Furthermore, if $X_{i}, i = 1, 2, \ldots, n$, with $n$ a positive integer, are pairwise independent random variables on $S$, then $V(X_{1}+X_{2}+ \cdots +X_{n}) = V(X_{1}) + V(X_{2}) + \cdots + V(X_{n})$
  \end{theorem} 
\end{frame}

\begin{frame}
  \frametitle{Bienayme's Formula Bevis}
  Vi ved at $V(X+Y) = E((X+Y)^{2}) - E(X+Y)^{2}$ Dermed:

  \begin{equation*}
    \begin{split}
      V(X+Y) &= E(X^{2} + 2XY + Y^{2}+Y^{2})-(E(X)+E(Y))^{2}\\
             &= E(X^{2}) +2E(XY)+E(Y^{2})-E(X)^{2}-2E(X)E(Y)-E(Y)^{2}
    \end{split}
  \end{equation*}

  Fordi X og Y er uafhængige, ved vi at $E(XY) = E(X)E(Y)$, og dermed:
  \begin{equation*}
    \begin{split}
      V(X+Y) &= (E(X^{2})-E(X)^{2})+E(Y^{2})-E(Y)^{2})\\
             &= V(X)+V(Y)\\
    \end{split}
  \end{equation*}
\end{frame}

\begin{frame}
  \frametitle{Chebyshev's Inequality}
  \begin{theorem}
    Let $X$ be a random variable on a sample space $S$ with proabbility functino $p$. If $r$ is a positive real number, then
    \begin{equation*}
      p(|X(s) - E(X)| \geq r) \leq V(X)/r^{2}
    \end{equation*}

  \end{theorem}
\end{frame}

\begin{frame}
  \frametitle{Chebyshev's Inequality Bevis}
  Lad $A$ være hændelsen:
  \begin{equation}
    A = \{s \in S | |X(S) - E(X)| \geq r\}
  \end{equation}
  Vi vil gerne bevise at $p(A) \leq V(X)/r^{2}$.

  \begin{equation}
    \begin{split}
      V(X) &= \sum_{s \in S} (X(s) - E(X))^{2}p(s)\\
           &= \sum_{s \in A} (X(s) - E(X))^{2}p(s) + \sum_{s \notin A} (X(s) - E(X))^{2} p(s)\\
      V(X) &\geq \sum_{s \in A} (X(s) - E(X))^{2}p(s) \\
           &= \sum_{s \in A} r^{2}p(s) \\
           &= r^{2}p(A)\\
      \frac{V(X)}{r^{2}} &\geq p(A)\\
    \end{split}
  \end{equation}
\end{frame}

\section{Chernoff Bounds}
\label{sec:label}

\begin{frame}
  \frametitle{Chernoff Bounds}
  
  \begin{itemize}
  \item Chernoff Bounds er bounds hvorpå du kan få et meget mere præcist bound end ved f.eks. Chebyshev's Inequality.
  \end{itemize}
\end{frame}

\begin{frame}[allowframebreaks]
  \frametitle{Theorem 13.42}
  Let $X_{1}+ X_2 + \cdots + X_{n}$, where $X_{i}$ takes the value 1 with the probability $p_{i}$, and the value 0 otherwise. $X = X_{1}+X_{2}+ \cdots + X_{n}$. Assume that $\mu \geq E[X]$. Then, for any $\delta > 0$, we have
  \begin{equation}
    Pr[X > (1+ \delta) \mu ] < \left[ \frac{e^{\delta}}{(1+\delta)^{(1 + \delta)}} \right]^{\delta}.
  \end{equation}
\end{frame}

\begin{frame}
  \frametitle{Chernoff Bounds Bevis}
  Læg først mærke til at for enhver $t > 0, p(X > (1 + \delta ) \mu) = p(e^{tX}> e^{t(1+\delta) \mu})$, da funktionen $f(x) = e^{tx}$ er monoton i $x$.

  Som det næste bruger vi nogle egenskaber for expectation. For et random variable, $Y$, har vi $\gamma p(Y > \gamma) \leq E[Y]$, af definition. Hvis vi sætter dette sammen, får vi:
  \begin{equation}
p(X > (1 + \delta ) \mu ) = p \left[ e^{tX} > e^{t(1 + \delta) \mu} \right] \leq e^{-t(1 + \delta) \mu } E \left[ e^{tX} \right]
\end{equation}

Som det næste skal vi bounde $E[e^{tX}]$. $E[e^{tX}] = E \left[ e^{t \sum_{i} X_{i}} \right] = E [ \Pi_{i} e^{tX_{i}}]$. For uafhængige variabler Y og Z, $E[YZ] = E[Y]E[Z]$. Dermed: $E[\pi_{i} e^{tX_{i}}] = \pi_{i} E[e^{tX_{i}}]$.

\end{frame}

\begin{frame}
  \frametitle{Chernoff Bevis part 2}
  
$e^{tX_{i}}$ er $e^{t}$ med sandsynligheden $p_{i}$, og ellers $e^{0}$. Dermed kan vi bounde dens expectation med:
\begin{equation}
  \label{}
E \left[ e^{tX_{i}} \right] = p_{i}e^{t}+ (1-p_{i})=1+p_{i}(e^{t}-1) \leq e^{p_{i}(e^{t}-1},
\end{equation}
hvor den sidste ulighed kommer fra at $1 + \alpha \leq e^{\alpha}$, så længe at $\alpha \geq 0$.
Hvis vi kombinerer ulighederne får vi:

\begin{equation}
  \label{}
p(X > (1 + \delta) \mu) \leq e^{-t(1 + \delta) \mu} E \left[ e^{tX} \right] = e^{-t(1+\delta) \mu} \Pi_{i} E \left[ e^{tX_{i}} \right] \leq e^{-t(1 + \delta) \mu} \Pi_{i} e^{p_{i(e^{t}-1}} \leq e^{-t(1+\delta) \mu}e^{\mu(e^{t}-1}.
\end{equation}
Givet at $t = \ln(1+\delta)$ får vi bounde.
\end{frame}

\begin{frame}
  \frametitle{Theorem 13.43}
 \begin{theorem}
   Let $X, X_{1}, X_{2}, \ldots, X_{n}$ and $\mu$ be defined as above. Then for any $1 > \mu > 0$, we have
   \begin{equation}
     \label{}
p(X < (1- \delta) \mu ) < e^{- \frac{1}{2} \mu \delta^{2}}
   \end{equation}

 \end{theorem} 
\end{frame}

\section{Weekly Note 3}
\label{sec:label}

\subsection{Notes on Indicator Random Variables}
\label{subsec:label}

\begin{frame}
  \frametitle{Indicator Random Variables}
  Lad $S$ være sample space, lad $A$ være en hændelse i $S$. Lad $X_{A}$ være en random variable som er 1 hvis $A$ hænder, og 0 ellers. Dermed:
  \begin{equation}
    \label{}
X_{A}(s) = \begin{cases}
  1 & \text{If } s \in A \\
  0 & \text{If } s \notin A \\
\end{cases}
  \end{equation}

  Vi kalder $X_{a}$ indicator variable for $A$.

\end{frame}

\begin{frame}
  \frametitle{Theorem 1}
 \begin{theorem}
Let $A$ be an event in a sample spce $S$ and let $X$ be the indicator variable for $A$. Suppose the probability that $A$ occurs is $p (p(A) = p)$. Then expected value and variance of $X$ is given by $E(X) = p, V(X) = p(1-p)$
 \end{theorem} 
\end{frame}

\begin{frame}[allowframebreaks]
  \frametitle{Theorem 1 Bevis}
  Lad os starte med expected value:
  \begin{equation}
    \begin{split}
      E(X) &= 0 \cdot p(X=0) + 1 \cdot p(X=1) \\
           &= 0 + 1 \cdot p \\
           &= p \\
    \end{split}
  \end{equation}

  Ved variansen bruger vi formlen $V(X) = E(X^{2}) - (E(X))^{2}$. Dette er nemt at gøre her, hvor hver indicator variable $X$ $X(s) = X(s)^{2}$ for alle $s \in S$. Nu får vi:

  \begin{equation}
    \begin{split}
      V(X) &= E(X^{2}) + (E(X))^{2}\\
           &= E(X) - (E(X))^{2}\\
           &= E(X)(1- E(X))\\
           &= p(1-p)\\
    \end{split}
  \end{equation}
\end{frame}

\begin{frame}
  \frametitle{Theorem 2}
  
  \begin{theorem}
Consider $n$ independent Bernoulli trials each of which have probability of success equal to $p$. Let $X$ be the random variable that denotes the number of successes in these $n$ trials. Then $E(X) = np$ and $V(X) = np(1-p)$.
  \end{theorem}
\end{frame}

\begin{frame}
  \frametitle{Theorem 2 Bevis}
  Vi har allerede bevist dette, dog ikke ved brug af indicator random variables.

  Lad $X_{i}, i = 1, 2, \ldots, n$ være indicator random varibl for hændelsen at den $i$'e trial eer en success (så $p(X_{i} =1) = p$  for all $i = 1, 2, \ldots, n$). Så er $X = \sum_{i=1}^{{X}}X_{i}$. Ved linearity of expectation, får vi $E(X) = \sum_{i=1}^{n}E(X_{i}) = \sum_{i=1}^{n}p = np$. Ydermere, da $X_{i}$ og $X_{j}$ er uafhængige random variables, for alle distinkte $i,j$, har vi at $V(X) = \sum_{i=1}^{n} V(X_{i}) = \sum_{i=1}^{n}p(1-p)=np(1-p)$
\end{frame}

\subsection{Notes on the Probabilistic Method}
\label{subsec:label}

\begin{frame}
  \frametitle{Markov's Inequality}
  \begin{theorem}[Markov's Inequality]
    For any non-negative random variable $X$ on a sample space $S$,
    \begin{equation}
      p(X \geq t) \leq \frac{E(X)}{t}
    \end{equation}
  \end{theorem}
\end{frame}

\begin{frame}
  \frametitle{Markov's Bevis}
  $E(X) = \sum_{i \in X(S)}p(X=i)i$. Siden $X(S) > 0$ for alle $s \in S$, får vi
  \begin{equation}
E(X) \geq \sum_{i \in X(S), i \geq t}p(X=i)i \geq t \sum_{i \in X(S), i \geq t}p(X = i) = p(X \geq t)t.
  \end{equation}
\end{frame}

\begin{frame}
  \frametitle{First Moment Principle}
  \begin{theorem}
If $E(X) \leq t$, then $p(X \leq t) > 0$
  \end{theorem}
\end{frame}

\begin{frame}
  \frametitle{First Moment Principle Bevis}
  Antag at $p(X \leq t) = 0$ . Så har vi at:
  \begin{equation}
E(X) = \sum_{i \in X(S)}^{}p(X=i)i \geq \sum_{i \in X(S), i > t}^{} p(X=i)i > t \sum_{i \in X(S), i \geq t}^{}p (X=i) =t
\end{equation},
hvor vi har brugt at $p(X \leq t) = 0$

\end{frame}


\end{document}
%%% Local Variables:
%%% mode: latex
%%% TeX-engine: xetex
%%% TeX-command-extra-options: "-shell-escape"
%%% End:
